% PAGES: 67-71 (printed pages 51-55; ends before the "2.5 Linear solvers using
% the LU decomposition without pivoting" heading, which begins later on
% printed page 55 and is covered by sec02_c3_2.tex)

% This file continues Section 2.4 "LU decomposition without pivoting", begun
% before printed page 51 (PDF page < 67) in another agent's file. The
% material below picks up mid-bullet: the previous page ends with the bullet
% "Compute the $i$-th row of $U$ and the $i$-th column of $L$:", and the
% framed pseudocode box below is that bullet's content. Do not open or edit
% the prior file.

\begin{center}
\fbox{\begin{minipage}{0.45\linewidth}
\noindent
for $k=i:n$\\
\hspace*{1.5em}$U(i,k) = A(i,k)$\\
\hspace*{1.5em}$L(k,i) = A(k,i)/U(i,i)$\\
end
\end{minipage}}
\end{center}

\noindent $\bullet$ Update the $(n-i)\times(n-i)$ lower right part of $A$ as follows:
\begin{equation}
A(i+1:n,i+1:n) \;=\; A(i+1:n,i+1:n) \;-\; L(i+1:n,i)U(i,i+1:n),
\end{equation}
which can be written entry by entry as follows:

\begin{center}
\fbox{\begin{minipage}{0.5\linewidth}
\noindent
for $j=(i+1):n$\\
\hspace*{1.5em}for $k=(i+1):n$\\
\hspace*{3em}$A(j,k) = A(j,k) - L(j,i)U(i,k)$\\
\hspace*{1.5em}end\\
end
\end{minipage}}
\end{center}

Further clarification on the recursive part of the LU decomposition can be
found in the example below for a $4\times 4$ matrix.

\begin{examplex}
Let
\begin{equation}
A \;=\; \begin{pmatrix} 2 & -1 & 3 & 0 \\ -4 & 5 & -7 & -2 \\ -2 & 10 & -4 & -7 \\ 4 & -14 & 8 & 10 \end{pmatrix}.
\end{equation}

Let $L$ and $U$ be the LU factors of $A$. The entries of the first row of $U$
are given by (2.22), i.e., $U(1,k)=A(1,k)$ for $k=1:4$.
\[
U(1,1)=2; \quad U(1,2)=-1; \quad U(1,3)=3; \quad U(1,4)=0.
\]

The entries of the first column of $L$ are given by (2.24), i.e.,
$L(k,1) = A(k,1)/U(1,1)$ for $k=1:4$:
\[
L(1,1)=1; \quad L(2,1)=\frac{-4}{U(1,1)}=\frac{-4}{2}=-2; \quad L(3,1)=\frac{-2}{U(1,1)}=\frac{-2}{2}=-1;
\]
\[
L(4,1)=\frac{4}{U(1,1)}=\frac{4}{2}=2.
\]

Then, the current forms of $L$ and $U$ are
\begin{equation}
L=\begin{pmatrix} 1 & 0 & 0 & 0 \\ -2 & 1 & 0 & 0 \\ -1 & L(3,2) & 1 & 0 \\ 2 & L(4,2) & L(4,3) & 1 \end{pmatrix};
\qquad
U=\begin{pmatrix} 2 & -1 & 3 & 0 \\ 0 & U(2,2) & U(2,3) & U(2,4) \\ 0 & 0 & U(3,3) & U(3,4) \\ 0 & 0 & 0 & U(4,4) \end{pmatrix}.
\end{equation}

The updated form of the $3\times 3$ matrix $A(2:4,2:4)$ is computed using
(2.36), from $A(2:4,2:4)$ obtained from (2.38) and with $L$ and $U$ given by
(2.39), as follows:
\begin{align*}
A(2:4,2:4) &= A(2:4,2:4) - L(2:4,1)U(1,2:4) \\
&= \begin{pmatrix} 5 & -7 & -2 \\ 10 & -4 & -7 \\ -14 & 8 & 10 \end{pmatrix} - \begin{pmatrix} -2 \\ -1 \\ 2 \end{pmatrix}(-1\ \ 3\ \ 0) \\
&= \begin{pmatrix} 5 & -7 & -2 \\ 10 & -4 & -7 \\ -14 & 8 & 10 \end{pmatrix} - \begin{pmatrix} 2 & -6 & 0 \\ 1 & -3 & 0 \\ -2 & 6 & 0 \end{pmatrix} \\
&= \begin{pmatrix} 3 & -1 & -2 \\ 9 & -1 & -7 \\ -12 & 2 & 10 \end{pmatrix}.
\end{align*}

Thus, the updated form of the $3\times 3$ matrix $A(2:4,2:4)$ is
\begin{equation}
A(2:4,2:4) \;=\; \begin{pmatrix} 3 & -1 & -2 \\ 9 & -1 & -7 \\ -12 & 2 & 10 \end{pmatrix}.
\end{equation}

Then,
\[
L(2:4,2:4)U(2:4,2:4) \;=\; \begin{pmatrix} 3 & -1 & -2 \\ 9 & -1 & -7 \\ -12 & 2 & 10 \end{pmatrix},
\]
which can be written as
\[
\begin{pmatrix} 1 & 0 & 0 \\ L(3,2) & 1 & 0 \\ L(4,2) & L(4,3) & 1 \end{pmatrix}
\begin{pmatrix} U(2,2) & U(2,3) & U(2,4) \\ 0 & U(3,3) & U(3,4) \\ 0 & 0 & U(4,4) \end{pmatrix}
=
\begin{pmatrix} 3 & -1 & -2 \\ 9 & -1 & -7 \\ -12 & 2 & 10 \end{pmatrix}.
\]

The unknown entries from the second row of $U$ and from the second column of
$L$ can be computed from the $3\times 3$ matrix above as follows:
\[
U(2,2)=3; \quad U(2,3)=-1; \quad U(2,4)=-2;
\]
\[
L(2,2)=1; \quad L(3,2)=\frac{9}{U(2,2)}=\frac{9}{3}=3; \quad L(4,2)=\frac{-12}{U(2,2)}=\frac{-12}{3}=-4.
\]

Then, the current forms of $L$ and $U$ are
\begin{equation}
L=\begin{pmatrix} 1 & 0 & 0 & 0 \\ -2 & 1 & 0 & 0 \\ -1 & 3 & 1 & 0 \\ 2 & -4 & L(4,3) & 1 \end{pmatrix};
\qquad
U=\begin{pmatrix} 2 & -1 & 3 & 0 \\ 0 & 3 & -1 & -2 \\ 0 & 0 & U(3,3) & U(3,4) \\ 0 & 0 & 0 & U(4,4) \end{pmatrix}.
\end{equation}

The updated form of the $2\times 2$ matrix $A(3:4,3:4)$ is computed using
(2.37), from $A(3:4,3:4)$ obtained from (2.40) and with $L$ and $U$ given by
(2.41), as follows:
\begin{align*}
A(3:4,3:4) &= A(3:4,3:4) - L(3:4,2)U(2,3:4) \\
&= \begin{pmatrix} -1 & -7 \\ 2 & 10 \end{pmatrix} - \begin{pmatrix} 3 \\ -4 \end{pmatrix}(-1\ \ -2) \\
&= \begin{pmatrix} -1 & -7 \\ 2 & 10 \end{pmatrix} - \begin{pmatrix} -3 & -6 \\ 4 & 8 \end{pmatrix} \\
&= \begin{pmatrix} 2 & -1 \\ -2 & 2 \end{pmatrix}.
\end{align*}

Thus, the updated form of the $2\times 2$ matrix $A(3:4,3:4)$ is
\begin{equation}
A(3:4,3:4) \;=\; \begin{pmatrix} 2 & -1 \\ -2 & 2 \end{pmatrix}.
\end{equation}

Then,
\[
L(3:4,3)U(3,3:4) \;=\; \begin{pmatrix} 2 & -1 \\ -2 & 2 \end{pmatrix},
\]
which can be written as
\[
\begin{pmatrix} 1 & 0 \\ L(4,3) & 1 \end{pmatrix}
\begin{pmatrix} U(3,3) & U(3,4) \\ 0 & U(4,4) \end{pmatrix}
=
\begin{pmatrix} 2 & -1 \\ -2 & 2 \end{pmatrix}.
\]

The unknown entries from the third row of $U$ and from the third column of
$L$ can be computed from the $2\times 2$ matrix above as follows:
\[
U(3,3)=2; \quad U(3,4)=-1;
\]
\[
L(3,3)=1; \quad L(4,3)=\frac{-2}{U(3,3)}=\frac{-2}{2}=-1.
\]

Then, the current forms of $L$ and $U$ are
\begin{equation}
L=\begin{pmatrix} 1 & 0 & 0 & 0 \\ -2 & 1 & 0 & 0 \\ -1 & 3 & 1 & 0 \\ 2 & -4 & -1 & 1 \end{pmatrix};
\qquad
U=\begin{pmatrix} 2 & -1 & 3 & 0 \\ 0 & 3 & -1 & -2 \\ 0 & 0 & 2 & -1 \\ 0 & 0 & 0 & U(4,4) \end{pmatrix}.
\end{equation}

The updated form of $A(4,4)$, which is a number, is computed using (2.37),
from $A(4,4)$ obtained from (2.42) and with $L$ and $U$ given by (2.43), as
follows:
\[
A(4,4) \;=\; A(4,4) - L(4,3)U(3,4) \;=\; 2-(-1)\cdot(-1) \;=\; 1,
\]
which corresponds to
\[
L(4,4)=1; \quad U(4,4)=1.
\]

We conclude that the matrix $A$ has an LU decomposition with the following
$L$ and $U$ factors:
\[
L=\begin{pmatrix} 1 & 0 & 0 & 0 \\ -2 & 1 & 0 & 0 \\ -1 & 3 & 1 & 0 \\ 2 & -4 & -1 & 1 \end{pmatrix};
\qquad
U=\begin{pmatrix} 2 & -1 & 3 & 0 \\ 0 & 3 & -1 & -2 \\ 0 & 0 & 2 & -1 \\ 0 & 0 & 0 & 1 \end{pmatrix}.
\]
\end{examplex}

The pseudocode for the LU decomposition without pivoting, given as a
function call $[L,U]=\text{lu\_no\_pivoting}(A)$, can be found in
Table~2.5.

\begin{table}[H]
\centering
\caption{Pseudocode for LU decomposition without pivoting}
\fbox{\begin{minipage}{0.9\linewidth}
Function Call:\\
$[L,U] = \text{lu\_no\_pivoting}(A)$
\bigskip

Input:\\
$A$ = nonsingular matrix of size $n$ with LU decomposition
\bigskip

Output:\\
$L$ = lower triangular matrix with entries $1$ on main diagonal\\
$U$ = upper triangular matrix\\
such that $A = LU$
\bigskip

\noindent
for $i = 1:(n-1)$\\
\hspace*{1.5em}for $k=i:n$\\
\hspace*{3em}$U(i,k) = A(i,k)$; \hfill\textit{// compute row $i$ of $U$}\\
\hspace*{3em}$L(k,i) = A(k,i)/U(i,i)$; \hfill\textit{// compute column $i$ of $L$}\\
\hspace*{1.5em}end\\
\hspace*{1.5em}for $j = (i+1):n$\\
\hspace*{3em}for $k = (i+1):n$\\
\hspace*{4.5em}$A(j,k) = A(j,k) - L(j,i)U(i,k)$;\\
\hspace*{3em}end\\
\hspace*{1.5em}end\\
end\\
$L(n,n) = 1$; $U(n,n) = A(n,n)$
\end{minipage}}
\end{table}

\begin{lemma}
The operation count for the LU decomposition without pivoting of an
$n\times n$ nonsingular matrix is
\begin{equation}
\frac{2}{3}n^3 \;+\; O(n^2).
\end{equation}
\end{lemma}

\begin{proof}
At step $i$, the ``for'' loop

\begin{center}
\fbox{\begin{minipage}{0.45\linewidth}
\noindent
for $k=i:n$\\
\hspace*{1.5em}$U(i,k) = A(i,k)$\\
\hspace*{1.5em}$L(k,i) = A(k,i)/U(i,i)$\\
end
\end{minipage}}
\end{center}

\noindent to compute the $i$-th row of $U$ and the $i$-th column of $L$
requires $n-i+1$ operations.

Also at step $i$, the double ``for'' loop

\begin{center}
\fbox{\begin{minipage}{0.5\linewidth}
\noindent
for $j=(i+1):n$\\
\hspace*{1.5em}for $k=(i+1):n$\\
\hspace*{3em}$A(j,k) = A(j,k) - L(j,i)U(i,k)$\\
\hspace*{1.5em}end\\
end
\end{minipage}}
\end{center}

\noindent to update $A(i+1:n,i+1:n)$ requires
\[
\sum_{j=i+1}^{n}\sum_{k=i+1}^{n} 2 \;=\; 2(n-i)^2
\]
operations.

When accounting for the outside ``for'' loop ``for $i=1:(n-1)$'', see
Table~2.5, we obtain that the operation count for the LU decomposition
without pivoting is
\begin{equation}
\sum_{i=1}^{n-1} \left(2(n-i)^2+n-i+1\right).
\end{equation}

Recall that
\begin{equation}
\sum_{l=1}^{p} l \;=\; \frac{p(p+1)}{2} \quad \text{and} \quad \sum_{l=1}^{p} l^2 \;=\; \frac{p(p+1)(2p+1)}{6};
\end{equation}
see, e.g., section 10.3.1 from Stefanica \cite{36}.

Then, by letting $l=n-i$ in (2.45) and using (2.46) for $p=n-1$, we obtain
that
\begin{align*}
\sum_{i=1}^{n-1}\left(2(n-i)^2+n-i+1\right)
&= \sum_{l=1}^{n-1}\left(2l^2+l+1\right)
\;=\; 2\sum_{l=1}^{n-1}l^2 + \sum_{l=1}^{n-1}l + \sum_{l=1}^{n-1}1 \\
&= 2\,\frac{(n-1)n(2n-1)}{6} \;+\; \frac{(n-1)n}{2} \;+\; n-1 \\
&= \frac{2n^3}{3} - \frac{n^2}{2} + \frac{5n}{6} - 1 \\
&= \frac{2}{3}n^3 + O(n^2);
\end{align*}
see (10.79) in Section 10.2.3 for the last equality.

We conclude that the operation count for the LU decomposition without
pivoting is $\dfrac{2}{3}n^3+O(n^2)$.
\end{proof}
